\section{Free heights and finite comparisons}\label{sec:heights}

The fixed spin boundary becomes a free boundary after duality. We first
establish this correspondence, including nonneutral charges, and then
record the height limits that we use from \cite{C}. The finite comparisons
must keep track of independent additive constants and prescribed data
consisting only of height differences. We spell out how the cited
inequalities apply with these conventions.

\subsection{Partition sums and fixed-spin duality}

Choose one orientation of every edge of a finite nearest-neighbor graph
\(G=(V,E)\), and write \(\nabla_eh=h_y-h_x\) for \(e=(x,y)\).
We call its vertices and edges primary, to distinguish them from the
auxiliary subdivisions used below.
An \emph{integral connection} is an antisymmetric edge function
\(\alpha\) with values in \(2\pi\Z\). The height variables lie in
\(2\pi\Z\), and their edge weights are
\begin{equation}\label{heights:weights}
 p_b(j)=e^{-b}I_j(b),\qquad
 \prod_{e\in E}p_b\left(\frac{\nabla_eh+\alpha_e}{2\pi}\right).
\end{equation}
Here \(I_j\) is the modified Bessel function. The symmetry
\(p_b(j)=p_b(-j)\) makes the chosen orientations immaterial. We will
occasionally include an observation penalty
\(\exp\{-\|Bh-v\|^2/2\}\), where \(B\) is a finite real linear map
and \(v\) is a prescribed real vector. Retained observation penalties
are sums of component-local forms in every cut topology considered:
each row of \(B\) is supported in one component. Within each component
they either annihilate its constant or confine it.

A component with no absolute constraint or confining observation is
summed modulo one common translation by \(2\pi\Z\). On a component whose
constant is confined by observations, all translation copies are summed.
An absolute height pin fixes that constant instead. These conventions
also apply after an edge cut produces several components. Thus fixing a
representative on each cut piece does not fix the relative translations
when the pieces are subsequently joined.

Write \(s\) for the collection of shifts, including \(\alpha\) and any
observation centers, and let \(Z_G^s\) be the corresponding partition
sum with counting measure and the preceding constant convention. The
superscript \(0\) means zero connection and zero observation centers,
with all cuts and quadratic penalties retained. Set
\begin{equation}\label{heights:sector-ratio}
                         F_G^s=\frac{Z_G^s}{Z_G^0}.
\end{equation}
If data \(u\) are prescribed on a set \(p\), then \(Z_G^s(u)\)
denotes the partition sum with those data imposed, not their probability.
Data can specify every height on \(p\), or every difference from a
chosen site of \(p\). In the second case the common height on \(p\)
is still free. Each such relative-data set lies in one component of
every topology being compared; data on several components are
specified separately. We prescribe only compatible lattice data.

Changing \(\alpha\) by the gradient of a \(2\pi\Z\)-valued function
and changing the height variable by the opposite function is an exact
reindexing of the sum. Observation centers and prescribed data must be
changed at the same time. We refer to this reindexing as an integral
gauge change. A connection is \emph{flat} on a region if its circulation
around every contractible cycle there is zero; on a multiply connected
region its remaining circulations are called its periods.

We use the classical Fourier-current duality in the formulation of
\cite[Section~2]{vEL}, with the fixed-spin boundary convention checked
in the following proof.

\begin{lemma}[Duality for the fixed spin boundary]\label{heights:duality}
Let \(\mathcal G_n\) have as vertices the centers of the lattice squares in
\([-1,1]^2\), and join centers of neighboring squares. This is the free
dual square grid. Given integers \(l_x\), indexed by the interior primal
vertices, choose an integral connection \(s\) on \(\mathcal G_n\) whose
circulation around \(x\) is \(2\pi l_x\), with a consistent choice of
orientation. Then
\[
       \E_n^0\exp\left(i\sum_xl_x\theta_x\right)=F_{\mathcal G_n}^s.
\]
Such a connection exists without a neutrality assumption. If holes
containing the charged vertices are removed, its restriction to the
remaining graph is flat, with periods specified by the enclosed charges.
On that graph the bare partition ratio depends only on those periods.
\end{lemma}

\begin{proof}
Edges lying along the spin boundary have both endpoints fixed at zero,
so their interaction factors cancel from the expectation. Expand every
other interaction as
\[
 e^{b\cos(\theta_x-\theta_y)}
       =\sum_{j\in\Z}I_j(b)e^{ij(\theta_x-\theta_y)}.
\]
Integration imposes divergence conditions only at interior vertices.
For the numerator these are the prescribed integer charges, with their
sign determined by the current orientation; the denominator has zero
divergence. There is no condition at a boundary vertex. Consequently
one can construct a current carrying each charge along a path to the
boundary, whether or not the total charge vanishes.

Subtract that fixed current in the numerator and rotate the remaining
edge variables onto the dual edges. Every retained primal edge crosses
exactly one edge of \(\mathcal G_n\). Every cycle of \(\mathcal G_n\) is generated by the
elementary cycles around interior primal vertices, so the rotated
zero-divergence current is a closed integral form. Integration along
dual paths represents it as \(\nabla h/(2\pi)\), uniquely modulo one
common height translation. The subtracted current gives the connection
\(\alpha\). The factors \(e^b\) needed to replace \(I_j(b)\) by \(p_b(j)\)
cancel. This proves the identity and also its positivity.

After holes are removed, any contractible cycle encloses no remaining
charge, while a cycle winding around a hole has circulation equal to
the total charge in that hole. Two integral flat connections with the
same periods differ by an integral gradient: integrate their difference
from a fixed vertex, using vanishing circulation for path independence.
The last assertion follows by gauge change.
\end{proof}

\subsection{The cited height limits}

For a finite graph \(G\), let \(A_G\) be its unscaled, unit-conductance
Laplacian, so that
\[
 \langle g,A_Gg\rangle=\sum_{\{x,y\}\in E(G)}(g_x-g_y)^2.
\]
Pairings here use counting measure. On a free square of side \(m\) in
lattice units, set \(g_m(x)=x_1/m\). The coefficient of \cite{C} is
\[
       a(b)=\lim_{m\to\infty}
                 \Var\langle h,A_Gg_m\rangle.
\]
The existence of this limit and the limit theorems below are cited
results about the Bessel weights \eqref{heights:weights}.

\begin{theorem}[Critical height coefficient {\cite[Section~6]{C}}]
\label{heights:critical-input}
For the free-box spin mass threshold \(b_c\) of Section~\ref{sec:intro},
the Bessel height coefficient satisfies \(a(b_c)=8\pi\).
\end{theorem}

The parameter and normalization in this theorem agree with ours:
\cite{C} uses the weight \(e^{-b}I_j(b)\), height spacing \(2\pi\),
and one copy of each unoriented edge. Its threshold is defined by the
same free-square spin mass. We use this as an independent theorem of
\cite{C}; it is not a consequence of the low-temperature spin theorem
of \cite{B}. In the remainder of the paper,
\[
                     b=b_c,\qquad a=8\pi,\qquad
                     K_*=\frac{a}{4\pi^2}=\frac2\pi.
\]

The following limit statements use three types of height observations.
A box average means
\[
                  h_Q=\frac1{|Q_n|}\sum_{x\in Q_n}h_x,
\]
where \(Q_n\) is the set of primary vertices in the discretized box.
A smear \(\langle h,f_n\rangle=\sum_xh_xf_n(x)\) on a free component
is defined only when \(\sum_xf_n(x)=0\) exactly. A continuum
mean-zero density \(f\) is represented by such profiles \(f_n\)
whose piecewise constant scaled densities \(n^2f_n\) converge to
\(f\) in \(L^2\); sampled profiles can be centered by subtracting
their discrete average.
For a free component with vertex set \(V_n\), its fractional mean is
\((|V_n|^{-1}\sum_{x\in V_n}h_x)\bmod 2\pi\), which is well-defined
on the height quotient.

\begin{theorem}[Free rectangle limits {\cite[Theorem~2.4]{C}}]
\label{heights:free-input}
Let \(U_n\) be free axis-parallel rectangles at mesh \(n^{-1}\), with
endpoints converging to those of a nondegenerate rectangle \(U\).
For any finite collection of smooth functions defined near
\(\overline U\), the observations
\(\langle h,A_{U_n}g\rangle\) converge jointly, in Laplace transforms
and all moments, to centered Gaussians with covariance
\[
                  a\int_U\nabla g\cdot\nabla g'.
\]
For bounded piecewise continuous mean-zero densities and the exactly
centered profiles just specified, height smears converge jointly in
Laplace transforms and all moments, with covariance
\(a(-\Delta_{U,N})^{-1}\), where the inverse is on mean-zero functions.
\end{theorem}

We next specify the domains for the mixed version of this theorem.
This will let us use rectangular partitions and cut regions without
requiring a limit theorem for arbitrary varying domains.

\begin{definition}[Fixed mixed geometry]\label{heights:geometry}
A fixed mixed geometry is a finite union of bounded axis-parallel
rectangular pieces with nondegenerate widths and faces. Pieces belonging
to one component meet along full interfaces, which may be subdivided
into segments. Components are polygonal Lipschitz domains, and no two
parts are joined solely through a point. A pin region is an aligned
finite union of positive-area boxes; pinning it means setting every
primary lattice height there to zero.

Discretizations use the nearest-neighbor graph at mesh \(n^{-1}\).
Boundaries may be rounded by a bounded number of lattice steps,
preserving the faces, corners, and positive-width passages. Observation
boxes lie within their respective components. A translated grid, such
as the grid of dual square centers, uses the same relative alignment
and rounding conventions.
\end{definition}

\begin{theorem}[Mixed height limits {\cite[Theorem~3.1]{C}}]
\label{heights:mixed-input}
In any fixed mixed geometry, the centered massless Bessel height law
with primary hard pins has Gaussian Laplace and moment limits for each
finite collection of bounded piecewise continuous density smears.
On a component \(U\) meeting a pin \(p\), the covariance is the inverse
of the form
\[
  a^{-1}\int_U|\nabla u|^2
  \quad\hbox{on}\quad
  \{u\in H^1(U):u=0\text{ almost everywhere on }p\}.
\]
Free walls have the variational Neumann condition. On an unpinned
component, the assertion concerns mean-free smears; the fractional mean
modulo \(2\pi\) tends to an independent uniform variable.

A fixed finite-rank nonnegative quadratic penalty in box averages gives
the corresponding Gaussian limit. When the penalty controls every
unpinned component constant, lattice translation copies are summed and
the limiting constants are integrated. The assertions hold along
every admissible sequence of rounded discretizations.
\end{theorem}

The next input concerns an interaction of rare pin events, not their
individual probabilities. It is this ratio that will control the cost
of restoring cut edges.

\begin{theorem}[Separated primary pins {\cite[Theorem~3.2]{C}}]
\label{heights:pin-input}
In a fixed mixed geometry, start with no hard pins and add the centered
penalty \(\|Bh\|^2/2\), where \(B\) is a fixed finite list of
component-local box averages controlling every component constant.
For positively separated pin regions \(H,p\), put
\[
 D_n(H,p)=
 \log\frac{\Pbb(h_H=0\mid h_p=0)}{\Pbb(h_H=0)}.
\]
Then \(D_n(H,p)\ge0\) and it converges to the Gaussian pin interaction
for the Hilbert-space inner product associated with
\[
         \mathcal T(u)=a^{-1}\sum_U\int_U|\nabla u|^2+\|Bu\|^2,
         \qquad u\in\bigoplus_U H^1(U).
\]
More precisely, with orthogonality and projections taken for
\(\mathcal T\),
\begin{equation}
 \begin{gathered}
 V_Q=\{u:u=0\text{ almost everywhere on }Q\}^{\perp_{\mathcal T}},
 \qquad \Pi_Q=\operatorname{proj}_{V_Q},\qquad
 C=\Pi_H|_{V_p},\\
             D_n(H,p)\longrightarrow
                  -\frac12\log\det(1-C^*C).
 \end{gathered}
 \tag{1}\label{eq:source-1}
\end{equation}
The operator \(C\) is Hilbert--Schmidt and has norm strictly below one.
\end{theorem}

The determinant formula is the Gaussian pin formula used in
\cite[Theorem~3.2]{C}; its mixed-domain formulation is given in
\cite[Proposition~3.14]{B}. Each use below fixes the geometry, the
observation boxes, and the positive separation before refining the
lattice. Further limits in cell size or geometric separation are
taken afterward. No conclusion uniform over all microscopic pin data
is included in these inputs.

\subsection{Affine comparisons and edge restoration}

We now pass from these limit inputs to comparisons on individual finite
graphs. The affine inequality of \cite[Lemma~2.2]{C} already includes
real observation centers and real numerator tilts. We recall its
Gaussian-chain proof before spelling out its application to cut graphs,
relative data, and independently released constants. All comparisons
in this subsection take place on fixed finite graphs.

\begin{lemma}[Finite comparison rules]
\label{heights:finite-comparisons}
Allow nonnegative quadratic observation penalties and compatible
linear equality constraints on a finite primary height graph.
The following comparisons hold.
\begin{enumerate}[label=\textup{(\roman*)}]
\item Increasing centered precision decreases centered covariance.
Consequently, in a proper centered law, the probability of an additional
centered linear equality increases under a centered precision increase.
The same assertion holds in a quotient law for equalities invariant
under its released constants.
\item A shifted-to-centered partition ratio is at most one and
decreases when the same nonnegative precision is added in the affine
and centered coordinates. Shifts may include integral connections,
nonzero prescribed data, and arbitrary real observation centers.
The monotonicity persists with a fixed real linear statistic in the
numerator exponent; the bound by one is not asserted for that
tilted ratio.
\item Equality constraints are obtained by increasing their quadratic
penalties to infinity. Free constants can be released when all data
and real tests in question are invariant under those constants.
\end{enumerate}
\end{lemma}

\begin{proof}
We first give the Gaussian comparison that will implement an edge
restoration. Let \(\Lambda\) be a lattice in a real Euclidean space,
let \(z+\Lambda\) be a translate, and let \(T\) be positive definite.
For \(Q\ge0\), \(t\ge0\), and a fixed real linear functional \(\ell\),
put
\[
 R(t)=
 \frac{\displaystyle\sum_{x\in z+\Lambda}
     e^{-\langle x,(T+tQ)x\rangle/2+\ell(x)}}
      {\displaystyle\sum_{x\in\Lambda}
     e^{-\langle x,(T+tQ)x\rangle/2}}.
\]
The relevant lattice-Gaussian inequalities originate in the work of
Regev and Stephens-Davidowitz \cite[Theorem~2.1 and Corollary~3.2]{RSD}.
The covariance under the numerator law dominates that under the
centered denominator law; see \cite[Lemma~2.1]{B}. If \(m_{z,\ell}\)
and \(\Cov_{z,\ell}\) denote the numerator mean and covariance,
and \(\Cov_0\) the denominator covariance, differentiation gives
\[
       \frac{d}{dt}\log R(t)
       =-\frac12\Tr Q\left(
          \Cov_{z,\ell}-\Cov_0+
          m_{z,\ell}m_{z,\ell}^{\mathsf T}\right)\le0.
\]
For a proper affine subspace, first complete the square in its
parallel span. The remaining offset is orthogonal for the current
precision \(T+tQ\); the covariance on the parallel span is that of a
translated, tilted lattice Gaussian with the restricted precision.
Thus its covariance still dominates the centered covariance, while
the raw mean \(m_{z,\ell}\) includes the fixed offset.
The centered maximum of a Gaussian lattice sum gives the untilted
bound \(R(t)\le1\).

To obtain the Bessel law, replace each primary edge by \(N\)
independent integer increments with normalized weight \(q_N^{j^2}\),
where \(q_N=b/(2N)\); their physical height increments are \(2\pi j\).
Uniformly for \(u\) in compact real sets,
\[
 \frac{\sum_{j\in\Z}q_N^{j^2}e^{uj}}{\sum_{j\in\Z}q_N^{j^2}}
   =1+\frac bN(\cosh u-1)+O_u(N^{-2}).
\]
Their sum tends to the probability weight \(p_b\), since its limiting
moment generating function is \(\exp\{b(\cosh u-1)\}\).
Matching endpoints, cycles, periods, and prescribed data gives
parallel affine lattices for the numerator and denominator.
Appending the residual coordinates \(Bh-v\) allows arbitrary real
observation centers. Thus adding a quadratic penalty or a matching
constraint is precisely the precision increase just considered.
Centered covariance decreases under such increases by the same
finite Gaussian comparisons; see \cite[Lemma~2.1]{B}.

All limits just described are on a fixed primary graph. A spanning
forest expresses every height difference as a sum of full-chain
increments, which have uniformly bounded exponential moments at
each fixed real argument. The other normalized edge weights are at
most one. This bounds the partition sums, their fixed real tilts,
and any polynomial factors. A confined component constant is
summed against its Gaussian penalty. Dominated convergence therefore
passes the comparisons to \(p_b\), as in \cite[Lemma~2.2]{C}.
Equality limits use monotone convergence of the penalties; only
nonempty affine fibers are used.
For centered pin monotonicity, first use a finite pin penalty.
The derivative of the logarithm of its expectation under an
additional centered precision \(Q\) is
\(\frac12\Tr Q(\Cov_{\rm unpinned}-\Cov_{\rm pinned})\ge0\).
Then increase the pin penalty to infinity.
The release of unconfined constants, including across cuts, is
justified next.
\end{proof}

\begin{lemma}[Cut comparison with prescribed data]\label{heights:cut-comparison}
Let \(G^{\rm cut}\) be obtained from \(G\) by deleting edges, with
the same observation penalties and compatible prescribed data on
both graphs. Each retained observation and each relative-data set
is component-local in \(G^{\rm cut}\). Whenever the free constants are left free by those
data,
\[
       \frac{Z_G^s(u)}{Z_G^0(0_p)}
       \le
       \frac{Z_{G^{\rm cut}}^s(u)}
            {Z_{G^{\rm cut}}^0(0_p)}.
\]
The conclusion also holds with absolute data fixing some constants
and with observations confining others. It holds for a fixed real
numerator tilt on affine edge increments if each deleted increment
is retained independently, with its tilt, on the cut side. Every
real test on a released component must annihilate its constant.
\end{lemma}

\begin{proof}
For a deleted edge \(e=(x,y)\), retain an independent chain with total
increment \(J_e\in2\pi\Z\). Its normalized sum without matching is
one. Restoring that edge imposes the equality
\[
                 J_e-h_y+h_x-\alpha_{xy}=0;
\]
the centered equality has \(\alpha_{xy}=0\).
On the enlarged affine coordinates this amounts to increasing the
coefficient of the square of the displayed residual to infinity.
The prescribed-data fibers in numerator and denominator are parallel
lattices. This remains true for relative data
\(h_z-h_{z_0}=u_z\), since these are linear equality constraints.
Lemma~\ref{heights:finite-comparisons} gives the comparison.
Integration of the matched chain contributes
\(p_b((h_y-h_x+\alpha_{xy})/(2\pi))\), and all independent-chain
normalizers cancel. For a real edge tilt, use \(J_e\) on the enlarged
graph. On matching it is precisely the affine edge increment in
the restored graph, so the tilted comparison applies with no
additional scalar.

Here is the promised treatment of constants. Before comparing, add
an auxiliary site penalty \(\delta\sum_xh_x^2/2\) and sum all
translations. For a free component with \(m\) vertices, a translation
orbit represented by \(h\), with mean \(\bar h\), contributes
\[
 e^{-\frac\delta2\sum_x(h_x-\bar h)^2}
 \sum_{k\in\Z}e^{-\frac{m\delta}{2}(\bar h+2\pi k)^2}.
\]
Divide its translation sum by
\[
                    c_{m,\delta}
                    =\frac1{2\pi}\sqrt{\frac{2\pi}{m\delta}}.
\]
Poisson summation shows that the quotient of the sum by
\(c_{m,\delta}\) tends uniformly in \(\bar h\) to one and is
uniformly bounded for small \(\delta\). The first factor is at
most one and tends pointwise to one. Forest exponential-moment
bounds again give domination, also for fixed real tilts invariant
under the released constant. Thus the regularized sum divided by
one \(c_{m,\delta}\) per free component tends to the quotient sum.
Within each shifted-to-centered ratio, its component factors cancel.
They cancel separately on the cut and restored sides, although
those graphs can have different numbers of components.

On a component already fixed by absolute data or confined by
observations, the auxiliary penalty is removed by ordinary dominated
convergence, with no divergent factor. This also covers comparison
of an unconfined ratio with a ratio in which new observations confine
every component: cancellation is performed separately in each ratio.
Semidefinite intermediate forms cause no difficulty, since the
auxiliary penalty makes every sum proper before the limit. Finally
take the chain limit at fixed primary graph as in
Lemma~\ref{heights:finite-comparisons}.
\end{proof}

In particular, the bare ratios satisfy
\begin{equation}
                 0<F_G^s\le1,\qquad F_G^s\le F_{G^{\rm cut}}^s,
 \tag{2}\label{eq:source-2}
\end{equation}
where the connection on the right is the restriction to the cut
graph. Deleting an additional region gives the same domain
comparison: first cut all its incident edges, then discard its
partition ratio, which is at most one.

\begin{corollary}[Restoration at prescribed data]\label{heights:restoration}
For an edge restoration with prescribed absolute or relative data
on \(p\), define
\[
       T^s(u)=\frac{Z_{\rm restored}^s(u)}{Z_{\rm cut}^s(u)},
       \qquad d=T^0(0_p).
\]
Use quotient constants on free components and sum all relative
integer translations when components are joined. Then
\[
                            0<T^s(u)\le d.
\]
\end{corollary}

\begin{proof}
The cut comparison reads
\[
 \frac{Z_{\rm restored}^s(u)}{Z_{\rm restored}^0(0_p)}
 \le
 \frac{Z_{\rm cut}^s(u)}{Z_{\rm cut}^0(0_p)}.
\]
Rearranging gives the claim. Every factor is positive and finite:
on a finite quotient graph the forest bound proves finiteness, and
all Bessel weights are positive. A joined component has only one
global quotient, so the numerator sums the relative constants that
the cut denominator quotients separately.
\end{proof}

The comparison is pointwise in the prescribed data. Its role is to
bound a restoration factor by its centered zero-data value. Lower
bounds will instead come from averaging separated pin comparisons;
they will not require a conditional limit uniform over microscopic
height data.
